% Minimal IEEEtran journal skeleton — the IEEE counterpart to templatePRIME.tex.
% Reference only; not part of the build. See README.md ("Switching to IEEEtran")
% for the preamble delta that turns paper-main.tex into an IEEE paper.
%
% IEEEtran.cls ships with TeX Live/MiKTeX (package `ieeetran`) and with the CI
% container, so nothing is vendored here. Only vendor a copy at the repo root if
% a venue pins a class version different from your distribution's.
%
% Class options: journal | conference | technote (+ lettersize | a4paper,
% draftcls for double-spaced review copy, onecolumn). IEEE Access uses its own
% option set — check the venue's author kit before submitting.
% `lettersize` comes from IEEE's current bare_jrnl sample but is not declared by
% IEEEtran v1.8b, so the log carries a benign "Unused global option [lettersize]".
\documentclass[lettersize,journal]{IEEEtran}

\usepackage[utf8]{inputenc} % allow utf-8 input
\usepackage[T1]{fontenc}    % use 8-bit T1 fonts
\usepackage[hypertexnames=false]{hyperref} % hyperlinks; hypertexnames=false avoids IEEEtran's duplicate float anchors
\usepackage{url}            % simple URL typesetting
\usepackage{booktabs}       % professional-quality tables
\usepackage{array}          % >{...} column modifiers
\usepackage{amsmath,amsfonts}
\usepackage{microtype}      % microtypography
\usepackage{graphicx}
\graphicspath{{media/}}     % organize your images and other figures under media/ folder
\usepackage[caption=false,font=footnotesize]{subfig} % subfigures — IEEE's recommended package (not subcaption)
\usepackage{stfloats}       % allow double-column floats at the bottom of a page
\usepackage{lipsum}         % filler text — delete in a real paper

\hyphenation{op-tical net-works semi-conduc-tor}

\begin{document}

%% Title. \thanks{} footnotes carry funding, manuscript dates, and the
%% corresponding-author note — IEEE has no separate \and-style author block.
\title{A Template for IEEE Transactions Style
\thanks{This work was supported by \dots}
\thanks{Manuscript received \dots; revised \dots.}}

\author{Author~One,~\IEEEmembership{Student Member,~IEEE,}
        Author~Two,~\IEEEmembership{Member,~IEEE,}
        and~Author~Three,~\IEEEmembership{Senior Member,~IEEE}% <-this % stops a space
\thanks{The authors are with the Department of \dots, University, City, Country
(e-mail: author@example.edu).}}

% Running heads: left = journal/volume line, right = short author + title.
\markboth{Journal of \LaTeX\ Class Files,~Vol.~14, No.~8, August~2021}%
{Author \MakeLowercase{\textit{et al.}}: A Template for IEEE Transactions Style}

\maketitle

\begin{abstract}
\lipsum[1]
\end{abstract}

\begin{IEEEkeywords}
First keyword, second keyword, more.
\end{IEEEkeywords}

% \IEEEPARstart{}{} sets the drop cap IEEE expects on the first paragraph of
% the introduction — and only there.
\section{Introduction}
\IEEEPARstart{T}{his} is the first paragraph. \lipsum[2]

\section{Method}
\label{sec:method}
\lipsum[3] See Section~\ref{sec:method}.

\subsection{Second-level heading}
\lipsum[4]
\begin{equation}
  y = \sum_{i=1}^{N} a_i x_i + b.
  \label{eq:example}
\end{equation}

\section{Results}
\lipsum[5] See Fig.~\ref{fig:fig1} and Table~\ref{tab:table}. IEEE cites by
bracketed number~\cite{shell2015ieeetran}, so \verb+\citet+-style inline forms
(available in the PRIME/arXiv template via natbib) do not apply here.

% Single-column float. Use figure* / table* to span both columns; with
% stfloats loaded, [!b] also works for double-column floats.
\begin{figure}[!t]
  \centering
  \fbox{\rule[-.5cm]{4cm}{3cm}}
  \caption{Sample figure caption.}
  \label{fig:fig1}
\end{figure}

% IEEE puts the table caption above the table and small-caps the title.
\begin{table}[!t]
  \caption{Sample Table Title}
  \label{tab:table}
  \centering
  \begin{tabular}{lll}
    \toprule
    Name     & Description     & Size ($\mu$m) \\
    \midrule
    Dendrite & Input terminal  & $\sim$100     \\
    Axon     & Output terminal & $\sim$10      \\
    Soma     & Cell body       & up to $10^6$  \\
    \bottomrule
  \end{tabular}
\end{table}

\section{Conclusion}
Your conclusion here.

\section*{Acknowledgments}
This work was supported in part by \dots

% References, two ways.
% (a) BibTeX — what paper-main.tex uses. IEEEtran.bst ships with the class:
%     \bibliographystyle{IEEEtran}
%     \bibliography{references}
%     Switch to these two lines once references.bib holds at least one cited
%     entry; against an empty .bib, IEEEtran's thebibliography errors out
%     ("perhaps a missing \item"), which is why this reference file uses (b).
% (b) A hand-written list, as IEEE's own bare_jrnl sample does:
\begin{thebibliography}{1}
\bibitem{shell2015ieeetran}
M.~Shell, ``How to use the IEEEtran \LaTeX\ class,'' IEEEtran class
documentation, v1.8b, Aug.\ 2015. [Online]. Available:
\url{https://ctan.org/pkg/ieeetran}
\end{thebibliography}

% Optional author biographies (IEEE journals only; drop for conference papers).
% \begin{IEEEbiography}[{\includegraphics[width=1in,height=1.25in,clip,keepaspectratio]{photo}}]{Author One}
% Biography text here.
% \end{IEEEbiography}

\end{document}
